% Options for packages loaded elsewhere
\PassOptionsToPackage{unicode}{hyperref}
\PassOptionsToPackage{hyphens}{url}
\PassOptionsToPackage{dvipsnames,svgnames,x11names}{xcolor}
%
\documentclass[
  10pt,
  a4paper,
]{article}
\usepackage{amsmath,amssymb}
\usepackage{setspace}
\usepackage{iftex}
\ifPDFTeX
  \usepackage[T1]{fontenc}
  \usepackage[utf8]{inputenc}
  \usepackage{textcomp} % provide euro and other symbols
\else % if luatex or xetex
  \usepackage{unicode-math} % this also loads fontspec
  \defaultfontfeatures{Scale=MatchLowercase}
  \defaultfontfeatures[\rmfamily]{Ligatures=TeX,Scale=1}
\fi
\usepackage{lmodern}
\ifPDFTeX\else
  % xetex/luatex font selection
  \setmainfont[]{Liberation Serif}
  \setsansfont[]{Liberation Sans}
  \setmonofont[]{DejaVu Sans Mono}
\fi
% Use upquote if available, for straight quotes in verbatim environments
\IfFileExists{upquote.sty}{\usepackage{upquote}}{}
\IfFileExists{microtype.sty}{% use microtype if available
  \usepackage[]{microtype}
  \UseMicrotypeSet[protrusion]{basicmath} % disable protrusion for tt fonts
}{}
\makeatletter
\@ifundefined{KOMAClassName}{% if non-KOMA class
  \IfFileExists{parskip.sty}{%
    \usepackage{parskip}
  }{% else
    \setlength{\parindent}{0pt}
    \setlength{\parskip}{6pt plus 2pt minus 1pt}}
}{% if KOMA class
  \KOMAoptions{parskip=half}}
\makeatother
\usepackage{xcolor}
\usepackage[margin=22mm]{geometry}
\setlength{\emergencystretch}{3em} % prevent overfull lines
\providecommand{\tightlist}{%
  \setlength{\itemsep}{0pt}\setlength{\parskip}{0pt}}
\setcounter{secnumdepth}{-\maxdimen} % remove section numbering
\ifLuaTeX
\usepackage[bidi=basic]{babel}
\else
\usepackage[bidi=default]{babel}
\fi
\babelprovide[main,import]{korean}
\ifPDFTeX
\else
\babelfont{rm}[]{Liberation Serif}
\fi
% get rid of language-specific shorthands (see #6817):
\let\LanguageShortHands\languageshorthands
\def\languageshorthands#1{}
\usepackage{xeCJK}
\setCJKmainfont{Noto Serif CJK KR}
\setCJKsansfont{Noto Sans CJK KR}
\setCJKmonofont{Noto Sans Mono CJK KR}
\xeCJKsetup{CJKspace=true}
\usepackage{xurl}
\usepackage{fvextra}
\usepackage{fancyhdr}
\AtBeginDocument{\DefineVerbatimEnvironment{verbatim}{Verbatim}{breaklines=true,breakanywhere=true,fontsize=\small}}
\setlength{\emergencystretch}{4em}
\setlength{\headheight}{15pt}
\pagestyle{fancy}\fancyhf{}
\fancyhead[L]{\small 중성자별 구조와 비방사 진동}
\fancyfoot[C]{\thepage}
\renewcommand{\headrulewidth}{0.3pt}
\setcounter{tocdepth}{2}
\ifLuaTeX
  \usepackage{selnolig}  % disable illegal ligatures
\fi
\usepackage{bookmark}
\IfFileExists{xurl.sty}{\usepackage{xurl}}{} % add URL line breaks if available
\urlstyle{same}
\hypersetup{
  pdftitle={A18 · 표의 취득 의미와 Hermite representation},
  pdflang={ko-KR},
  colorlinks=true,
  linkcolor={Maroon},
  filecolor={Maroon},
  citecolor={Blue},
  urlcolor={Blue},
  pdfcreator={LaTeX via pandoc}}

\title{A18 · 표의 취득 의미와 Hermite representation}
\usepackage{etoolbox}
\makeatletter
\providecommand{\subtitle}[1]{% add subtitle to \maketitle
  \apptocmd{\@title}{\par {\large #1 \par}}{}{}
}
\makeatother
\subtitle{과제 큐시트}
\author{}
\date{}

\begin{document}
\maketitle

\setstretch{1.12}
\textbf{선행:} A10, A13의 Hermite 유도　 \textbf{독립 작업:} 3--5시간
(계획치)

\subsection{표를 읽었다면 별에 넣을 EOS도 준비된
것일까?}\label{uxd45cuxb97c-uxc77duxc5c8uxb2e4uxba74-uxbcc4uxc5d0-uxb123uxc744-eosuxb3c4-uxc900uxbe44uxb41c-uxac83uxc77cuxae4c}

파일의 열과 단위는 읽었지만 보간 구간의 미분은 검사하지 않았다. 이
상태에서 성공이라고 부를 수 있는 작업과 아직 남은 작업을 나누자.

먼저 한두 문장으로 답한 뒤 아래 설명을 읽자. 첫 예측이 틀렸거나 아직
모르겠다고 적는 것은 감점하지 않는다.

\subsection{필요한 관계와 읽을
위치}\label{uxd544uxc694uxd55c-uxad00uxacc4uxc640-uxc77duxc744-uxc704uxce58}

하나의 \(P(\mu)\)에서 \(n=P'(\mu)\), \(\epsilon=\mu n-P\),
\(\chi=n/(\mu P''(\mu))\)를 함께 계산한다. 원표의 에너지 열이 이 관계를
만족하는지 확인하는 일과, 선택한 보간 함수가 구간 전체에서 허용되는지는
다른 검사다. 제9장의 유도 뒤에 A20에서 local-cell 평가를 직접 호출한다.

\begin{itemize}
\tightlist
\item
  \path|notes/chapters/ch09.md|
\item
  \path|methods/src/eos/ns_table_eos.f90|
\item
  \path|tools/normalize_compose.py|
\item
  \path|tools/inspect_table.py|
\item
  \path|examples/method_calls/table_query.f90|
\end{itemize}

\subsection{계산으로
옮기기}\label{uxacc4uxc0b0uxc73cuxb85c-uxc62euxae30uxae30}

제공 함수의 호출이 기본 경로다. 구간 보간·역함수·적응 적분·matching의
내부 코드를 작성하지 않고, 입력 조건과 반환된 근거를 해석한다.

\begin{enumerate}
\def\labelenumi{\arabic{enumi}.}
\tightlist
\item
  동봉 synthetic 표를 inspect\_table로 읽고 제공된 허용조건 검사를
  호출한다.
\item
  동봉 ID134 원자료와 metadata를 읽고 source order를 보존한다.
\item
  ID134의 196개 구간 중 거절된 72개를 확인한다. 이 표현을 임의로
  평활화하지 않는다.
\item
  동봉 ID69의 의미와 native 구간을 확인하고 A20의 제공 함수로 값과
  미분을 함께 읽는다.
\end{enumerate}

\newpage

\subsection{실행해서 무엇을 볼
것인가}\label{uxc2e4uxd589uxd574uxc11c-uxbb34uxc5c7uxc744-uxbcfc-uxac83uxc778uxac00}

패키지 루트에서 실행한다. \texttt{-\/-plan}은 경로만 확인하며 계산
성공을 뜻하지 않는다.

\begin{verbatim}
python tools/run_assignment.py A18 --run \
  --output ../A18_result.json
python tools/inspect_table.py \
  fixtures/tables/hermite_consistent.csv \
  --output ../synthetic_table
\end{verbatim}

ID134의 \path|matching_inconsistent|와 \texttt{model\_ready:\ false}는
허용조건의 거절이다. 종료코드 2를 성공한 별의 입력으로 해석하지 않는다.

\subsection{결과를
설명하기}\label{uxacb0uxacfcuxb97c-uxc124uxba85uxd558uxae30}

원자료의 의미와 선택한 함수의 허용조건을 따로 판단한다. 거절된 구간을
지우거나 평활화해서 같은 물리 문제라고 보고하지 않는다.

\textbf{남길 증거.} raw 열 정의와 SI 변환; admitted/rejected interval
설명; 각 표의 허용·거절 이유와 실제 별 계산에 사용할 함수의 정의

\textbf{판단 기준.} 2점 예: raw identity, 단위 의미, representation
admissibility를 별도 evidence로 제시한다. 최초 예측의 정오보다, 결과를
본 뒤 이유를 고치고 다른 입력에도 적용하는지를 본다.

\subsection{막혔을 때}\label{uxb9c9uxd614uxc744-uxb54c}

\textbf{힌트 1.} mu 비단조나 duplicate를 sort·삭제하여 문제를 숨기지
않는다.

\textbf{힌트 2.} 동일한 P,μ의 두 행이 값 중복인지 phase jump인지
의미부터 확인한다.

\subsection{다음 만남에 다시
묻기}\label{uxb2e4uxc74c-uxb9ccuxb0a8uxc5d0-uxb2e4uxc2dc-uxbb3buxae30}

다음 회차에는 중복 행의 의미를 묻는다. 표의 순서와 물질 상태를 확인하지
않고 정렬·삭제하면 무엇을 잃을 수 있는지 설명한다. 먼저 짧게 답한 뒤
본문과 자신의 기록을 열어 확인한다.

\textbf{범위.} Synthetic 예제의 전체 밀도 범위를 실제 별에 외삽하지
않는다. 별 계산은 명시한 ID69의 유효구간과 양의 표면압력에 한정한다.

\end{document}
